% c07 —— §11 轨形上同调与 Vafa 公式（p027–p028）、§12 环面簇前半（p029–p030）
% 编号公式 (11.1)–(11.15)、(12.1)–(12.23)，共 38 条
% 开篇承接 p026 末尾 (10.38) 之后的「で定義される。」

平凡群 $1$ 的转置为 $\check G_{\max}$，而 $\mathrm{SL}_n(\C)\cap G_{\max}$ 与 $\langle J\rangle$ 在转置下互换。

[CR10] 的主定理是拓扑镜像对称层面辛几何一侧的所谓 Calabi–Yau/Landau–Ginzburg 对应（Calabi–Yau/Landau–Ginzburg correspondence）；上述关于 Calabi–Yau 轨形的拓扑镜像对称，正是通过归结为 [Kra] 所证明的关于 Landau–Ginzburg 轨形的拓扑镜像对称而得到证明的。

% ================================================================ p027
\section{轨形上同调与 Vafa 公式}

光滑 Deligne–Mumford 栈的轨形上同调（orbifold cohomology）定义为惯性栈的上同调：
\begin{equation}
    H^{p,q}_{\mathrm{orb}}(\mathbb{Y}) := \bigoplus_{\cZ\subset\mathrm{I}\mathbb{Y}} H^{p-a(\cZ),q-a(\cZ)}(\cZ). \tag{11.1}
\end{equation}

这里，把惯性栈（inertia stack）定义为栈的范畴中的纤维积
\begin{equation}
    \mathrm{I}\mathbb{Y} := \mathbb{Y}\times_{\mathbb{Y}\times\mathbb{Y}}\mathbb{Y}, \tag{11.2}
\end{equation}
而对整体商，它具有
\begin{align}
    \mathrm{I}[M/G] &\cong \left[\left(\coprod_{g\in G}M^{g}\right)\middle/ G\right] \tag{11.3}\\
    &\cong \coprod_{g\in G/\operatorname{Ad}(G)} \left[M^{g}/Z(g)\right] \tag{11.4}
\end{align}
这样的具体描述。这里，$M^{g}$ 是 $g$ 作用的固定点集，$h\in G$ 通过
\begin{equation}
    M^{g}\to M^{hgh^{-1}},\quad x\mapsto hg \tag{11.5}
\end{equation}
作用于 $\coprod_{g\in G}M^{g}$。又，$Z(g)$ 是 $g$ 的中心化子。直观地说，「$\mathrm{I}\mathbb{Y}$ 的点」就是「$\mathbb{Y}$ 的点与该点自同构组成的对」。把轨形上同调定义为惯性栈的上同调、而非 $\mathbb{Y}$ 的上同调，原因在于：来自扭结稳定曲线（twisted stable map）的赋值映射（evaluation map）取值于 $\mathrm{I}\mathbb{Y}$（更准确地说，是刚性化惯性栈（rigidified inertia stack）），而非 $\mathbb{Y}$。

% [Tutor Note: 原文 (11.5) 右端作 x ↦ hg，按作用定义应为 x ↦ hx；此处按原文保留，请审校确认。]
% [Tutor Note: 原文此处作「捻れ安定曲線（twisted stable map）」，日文「曲線」（曲线）与括注英文 map 不一致，疑「曲線」为「写像」之误；译文按原文保留。]

(11.1) 右端的直和遍历惯性栈 $\mathrm{I}\mathbb{Y}$ 的连通分支 $\cZ$。各 $\cZ$ 仍然是光滑的，故其粗模概形 $Z$ 至多具有商奇点，其 $\Q$ 系数上同调带有纯（pure）Hodge 结构。

把给定连通分支 $\cZ\subset\mathrm{I}\mathbb{Y}$ 上的一般点用 $\mathbb{Y}$ 的局部表示 $[M/G]$ 以及 $x\in M$、$g\in G$ 表为 $(x,g)$ 时，把 $g$ 在 $T_xM$ 上作用的特征值之集合计重数写作 $\{\exp(2\pi\sqrt{-1}a_i^{g})\}_{i=1}^{\dim M}$（其中 $0\le a_i^{g}<1$），则 $\cZ$ 的年龄（age）定义为
\begin{equation}
    a(g) := \sum_{i=1}^{\dim M} a_i^{g}. \tag{11.6}
\end{equation}
(11.1) 中按年龄所作的次数移位（age shift），经由川崎彻郎的 Riemann–Roch 定理（Kawasaki–Riemann–Roch theorem）[Kaw79]，是使轨形（量子）杯积保持次数所必需的。轨形（量子）上同调的入门可参见例如 [Abr08]。

\begin{equation}
    \mathbb{Y} := \left[\left(f^{-1}(0)\setminus\mathbf{0}\right)\middle/\mathbb{G}_m\right]. \tag{11.7}
\end{equation}

% ================================================================ p028
其惯性栈由
\begin{equation}
    \mathrm{I}\mathbb{Y} = \coprod_{\gamma\in\mathbb{G}_m}\mathbb{Y}_{\gamma} \tag{11.8}
\end{equation}
给出。这里，对 $\gamma\in\mathbb{G}_m$，把 $f\colon\mathbb{A}^n\to\mathbb{A}^1$ 限制到 $\gamma$ 不变子空间 $(\mathbb{A}^n)^{\gamma}\subset\mathbb{A}^n$ 上所得的多项式记作 $f_{\gamma}$ 时，$\mathbb{Y}_{\gamma}$ 定义如下：
\begin{equation}
    \mathbb{Y}_{\gamma} := \left[\left(f_{\gamma}^{-1}(0)\setminus\mathbf{0}\right)\middle/\mathbb{G}_m\right]. \tag{11.9}
\end{equation}
因此，轨形上同调由
\begin{equation}
    H^{p,q}_{\mathrm{orb}}(\mathbb{Y}) \cong \bigoplus_{\gamma\in\mathbb{G}_m} H^{p-a(\gamma),q-a(\gamma)}(\mathbb{Y}_{\gamma}) \tag{11.10}
\end{equation}
给出。$(\mathbb{A}^n)^{\gamma}=\mathbf{0}$ 时 $\mathbb{Y}_{\gamma}$ 为空集，故 (11.10) 中的直和为有限和。

由 Griffiths 对射影超曲面上同调的计算 [Gri69]，经 Steenbrink 推广到加权射影超曲面 [Ste77]，$Y$ 的中间维数上同调由
\begin{equation}
    H^{p,n-2-p}(X) \cong \operatorname{Jac}(f)^{h(n-1-p)-\sum_{i=1}^{n}a_i} \tag{11.11}
\end{equation}
给出。这里，对右端的元素 $P$，左端的元素由
\begin{equation}
    \operatorname{Res}\frac{P\Omega}{f^{\,n-1-p}} \tag{11.12}
\end{equation}
定义。其中，$\Omega$ 由下式给出：
\begin{equation}
    \Omega = \sum_{i=1}^{n}(-1)^{i}x_i\bigwedge_{j\neq i}\dd x_j. \tag{11.13}
\end{equation}
另见 [Dol82, Theorem 4.3.2] 与 [CK99, Section 5.3]。当 Calabi–Yau 条件
\begin{equation}
    h = \sum_{i=1}^{n}a_i \tag{11.14}
\end{equation}
成立时，
\begin{equation}
    H^{p,n-2-p}(X) \cong \operatorname{Jac}(W)^{h(n-2-p)} \tag{11.15}
\end{equation}
成立。此外，在中间维数以外，由 Lefschetz 定理，上同调与加权射影空间的上同调一致。$\mathbb{Y}_{\gamma}$ 的上同调也可同样计算。

% [Tutor Note: 原文此段正文作「Y の…コホモロジー」，而 (11.11) 与 (11.15) 左端写作 X，记号不一致；此处按原文保留。]

Vafa 公式直接计数的并不是 $\mathbb{Y}$ 的轨形 Hodge 数，而是 Landau–Ginzburg 轨形 $([\mathbb{A}^n/G],f)$ 的 Hodge 数。其中令 $G := \mu_h := \operatorname{Ker}\bigl((-)^h\colon\mathbb{G}_m\to\mathbb{G}_m\bigr)$。Landau–Ginzburg 轨形及其 Hodge 数等概念此处不予定义，但其基本哲学是：把轨形与其上的全纯函数构成的对视为空间。全纯函数若为 Bott–Morse 型，则人们期待：考虑 Landau–Ginzburg 模型与考虑由临界点构成的子空间几乎是等价的。Landau–Ginzburg 轨形与 Calabi–Yau 超曲面之间的关系，就是第 10 节末尾提到的 Calabi–Yau/Landau–Ginzburg 对应。这也涉及关于 $G\subset\mathrm{SL}_n$ 的 McKay 对应。对可逆多项式，经由 Calabi–Yau/Landau–Ginzburg 对应可知，Vafa 公式正确地计算了 Calabi–Yau 超曲面的 Hodge 数。关于不具有可逆多项式的拟光滑 Calabi–Yau 超曲面的轨形 Hodge 数，可参见例如 [Bat]。

% ================================================================ p029
\section{环面簇}

带有代数群作用、且该作用具有开稠密轨道的流形，称为预齐性簇（prehomogeneous variety）。设 $X$ 为 $G$ 的预齐性簇，开稠密轨道记作 $U$。把 $G$ 在点 $x\in X$ 处的稳定化子记作
\begin{equation}
    Z_x := \{g\in G\mid t\cdot x = x\}, \tag{12.1}
\end{equation}
则
\begin{equation}
    Z_{g\cdot x} = \operatorname{Ad}_g(Z_x) := \{ghg^{-1}\in G\mid h\in Z_x\} \tag{12.2}
\end{equation}
成立。对任意 $x\in X$，过 $x$ 的 $G$ 作用轨道 $\cO_x := G\cdot x$ 是齐性空间，且
\begin{equation}
    \cO_x \cong G/Z_x \tag{12.3}
\end{equation}
成立。若 $G$ 可交换，则同一 $G$ 轨道上的各点具有相同的稳定化子；又，当 $x$ 含于 $\cO_y$ 的闭包中时 $Z_x\supset Z_y$。因此，开稠密轨道的稳定化子同时也是所有点的稳定化子，对它取商之后，可以假定 $G$ 自由地作用于开稠密轨道（从而 $G$ 对 $X$ 的作用是忠实的（faithful））。
% [Tutor Note: 原文作「開稠密軌道の固定化部分群は全ての点の固定化部分群にもなっており」，若按前句 Z_x⊃Z_y 的包含方向理解，开稠密轨道（的点）的稳定化子应包含于所有点的稳定化子之中；此处按原文照译，请审校确认。]

% [Tutor Note: 原文 (12.1) 的集合条件写作 t · x = x，而变元为 g ∈ G，按上下文应为 g · x = x；此处按原文保留。]

乘法群 $\mathbb{G}_m$ 的直积称为环面\footnote{一般而言，在平坦准紧（fpqc）拓扑下局部地与乘法群的直积同构的群概形，称为环面。该条件意味着：在域上，基变换到代数闭域后同构于乘法群的直积；在代数闭域上则就是乘法群的直积。}，而作为环面的预齐性簇的正规代数簇称为环面簇。

环面按子群取商仍是环面，故可以假定环面簇含有一个环面作为开稠密轨道。换言之，环面簇是指环面 $T$ 到正规簇 $X$ 的部分紧化，且 $T$ 在自身上的作用延拓为在 $X$ 上的作用。环面簇的典型例子有仿射空间 $\mathbb{A}^n$、射影空间 $\PP^n$，或者更一般地，加权射影空间 $\PP(a_1,\dots,a_{n+1})$ 等。在加权射影空间 $\PP(a_1,\dots,a_{n+1})$ 上，环面 $(\mathbb{G}_m)^n$ 通过 $(\alpha_1,\dots,\alpha_{n+1})\cdot[x_1:\cdots:x_{n+2}] := [\alpha_1x_1:\cdots:\alpha_{n+1}x_{n+1}:x_{n+2}]$ 作用，由所有坐标均不为 $0$ 的点构成的子集成为开稠密轨道。此外，环面自身也是环面簇的例子。

% [Tutor Note: 原文此处作 [x_1 : ⋯ : x_{n+2}] 与 x_{n+2}，但 P(a_1,…,a_{n+1}) 只有 n+1 个坐标，相应坐标应为 x_1,…,x_{n+1}；又作用被说成 (G_m)^n 而参数却写作 α_1,…,α_{n+1}。两处均按原文保留。]

给定环面 $T$ 时，令 $T$ 的单参数子群与特征标构成的自由 Abel 群为
\begin{gather}
    N := \Hom(\mathbb{G}_m,T),\quad M := \Hom(T,\mathbb{G}_m), \tag{12.4}\\
    N_{\R} := N\otimes\R,\quad M_{\R} := M\otimes\R. \tag{12.5}
\end{gather}
用 $M\times\R$ 的有限子集 $S$，由
\begin{equation}
    \left\{n\in N_{\R}\mid \text{对任意}\ (m,\lambda)\in S\ \text{均有}\ \langle n,m\rangle+\lambda\ge 0\right\} \tag{12.6}
\end{equation}
定义的 $N_{\R}$ 的子集称为多面体（polyhedron），其中有界者称为多面体（polytope）。后者也可定义为 $N_{\R}$ 的有限个点的凸包。特别地，作为 $N$ 的有限子集的凸包（convex hull，即包含该集合的最小凸集）给出的多面体称为格多面体（lattice polytope）。

关于 $\R^{>0}$ 作用不变的 $N_{\R}$ 的子集称为锥（cone），而为锥的多面体（polyhedron）称为多面锥（polyhedral cone）。$N_{\R}$ 的子集 $\sigma$ 为凸（convex）是指：对 $\sigma$ 的任意 2 点，连接它们的线段包含在 $\sigma$ 中。

% ================================================================ p030
多面锥在不含直线时称为强凸（strongly convex）。由 $N$ 的有限子集生成的（即作为包含该集合的最小锥给出的）多面锥称为有理多面锥（rational polyhedral cone）。对取面操作与取交操作封闭的有理强凸多面锥的集合 $\Sigma$ 称为扇（fan）。

把环面簇 $X$ 分解为 $T$ 轨道时，各个 $T$ 轨道都可用 $T$ 的某个单参数子群 $n\in N$ 表为
\begin{equation}
    n(0)(T) := \left\{\lim_{s\to 0}n(s)(t)\in X\mid t\in T\subset X\right\}. \tag{12.7}
\end{equation}
对 $T$ 轨道 $\cO\subset X$，把包含 $\{n\in N\mid n(0)(T)=\cO\}$ 的最小闭凸锥记作 $\sigma_{\cO}$，则
\begin{equation}
    \Sigma := \left\{\sigma_{\cO}\mid \cO\subset X\ \text{为}\ T\ \text{轨道}\right\} \tag{12.8}
\end{equation}
构成扇。

作为例子，先考虑 $X=T=(\mathbb{G}_m)^2$ 的情形。$X$ 只有一个 $T$ 轨道，对应的闭凸锥由
\begin{equation}
    \sigma_T := \{(0,0)\}\subset N_{\R}\cong\R^{2} \tag{12.9}
\end{equation}
给出，故与 $X$ 对应的扇为
\begin{equation}
    \Sigma = \{\sigma_T\}. \tag{12.10}
\end{equation}
其次，把 $T$ 作部分紧化得到 $\mathbb{A}^1\times\mathbb{G}_m$，则 $\mathbb{A}^1\times\mathbb{G}_m$ 上的 $T$ 作用 $(t_1,t_2):=(x_1,x_2)\mapsto(t_1x_1,t_2x_2)$ 的轨道为
% [Tutor Note: 原文把作用写作 (t_1,t_2):=(x_1,x_2)↦(t_1x_1,t_2x_2)，按上下文「:=」应为「·」（与 (12.1) 的 t·x 同类笔误）；按原文保留。]
\begin{align}
    &T = \{(x_1,x_2)\in\mathbb{A}^1\times\mathbb{G}_m\mid x_1\neq 0,\ x_2\neq 0\}, \tag{12.11}\\
    &\cO_1 = \{(x_1,x_2)\in\mathbb{A}^1\times\mathbb{G}_m\mid x_1=0,\ x_2\neq 0\}, \tag{12.12}\\
    &\tag{12.13}
\end{align}
这 2 个，对应的闭凸锥为
\begin{align}
    &\sigma_T = \{(0,0)\}, \tag{12.14}\\
    &\sigma_{\cO_1} = \{(n_1,0)\in N_{\R}\mid n_1\ge 0\}. \tag{12.15}
\end{align}
再把它部分紧化为 $\mathbb{A}^2$，则作用的轨道为
\begin{align}
    &T = \{(x_1,x_2)\in\mathbb{A}^2\mid x_1\neq 0,\ x_2\neq 0\}, \tag{12.16}\\
    &\cO_1 = \{(x_1,x_2)\in\mathbb{A}^2\mid x_1=0,\ x_2\neq 0\}, \tag{12.17}\\
    &\cO_2 = \{(x_1,x_2)\in\mathbb{A}^2\mid x_1\neq 0,\ x_2=0\}, \tag{12.18}\\
    &\cO_{12} = \{(0,0)\} \tag{12.19}
\end{align}
这 4 个，对应的闭凸锥为
\begin{align}
    &\sigma_T = \{(0,0)\}, \tag{12.20}\\
    &\sigma_{\cO_1} = \{(n_1,0)\in N_{\R}\mid n_1\ge 0\}, \tag{12.21}\\
    &\sigma_{\cO_2} = \{(0,n_2)\in N_{\R}\mid n_2\ge 0\}, \tag{12.22}\\
    &\sigma_{\cO_{12}} = \left\{(n_1,n_2)\in N_{\R}\mid n_1\ge 0\ \text{且}\ n_2\ge 0\right\}. \tag{12.23}
\end{align}

% [Tutor Note: 原文 (12.13) 为一个不承载公式内容的空编号行，此处按原版式保留为空白编号行。]

% CHUNK c07 DONE pages p027–p030
